\section{Model}
\label{sec:model}

The model tracks subjective sleepiness $s(t)$ on the KSS scale (1--9). It also tracks perceived workload on the raw NASA-TLX scale; that part is a simple, unevaluated feature and is described in Supplementary Section~S1. All parameters are listed in Table~\ref{tab:params}.

\subsection{Prior: the Three-Process Model}
For a person who woke at time $t_w$, let $\tau = t - t_w$ in hours and $h(t)$ the local clock hour. Alertness is the sum of a homeostatic, a circadian, an ultradian and a sleep-inertia component~\cite{akerstedt2004predictions,ingre2014validating}:
\begin{align}
  S(t) &= L + (S_w - L)\,e^{-\delta_w \tau}, \label{eq:S}\\
  C(t) &= M \cos\big(2\pi (h(t) - \phi)/24\big), \\
  U(t) &= U_m - U_a \cos\big(2\pi (h(t) - \phi_U)/12\big), \\
  W(t) &= W_0\, e^{-\delta_i \tau}, \\
  A(t) &= S(t) + C(t) + U(t) + W(t).
\end{align}
During sleep, $S$ recovers linearly at rate $\delta_s (H - B)$ while it is below the breakpoint $B = 12.2$ and exponentially towards $H$ above it~\cite{akerstedt2008accounting,ingre2014validating},
\begin{equation}
  \frac{dS}{dt} = \begin{cases} \delta_s (H - B), & S < B,\\ \delta_s (H - S), & S \geq B. \end{cases}
\end{equation}
The level at waking, $S_w$, is computed by running these equations over the recorded sleep episodes, preceded by ten habitual nights (7.5~h), so that a run of short nights builds up sleep debt. Alertness is mapped to the KSS and clipped,
\begin{equation}
  p(t) = \min\!\big(9, \max(1,\; a_K - b_K\,A(t))\big).
  \label{eq:prior}
\end{equation}
For a person who wakes at 07:00 after 8~h of sleep, the prior is KSS~7.0 at waking (sleep inertia), 3.7 at 09:00, 3.1 at 16:00, 5.7 at 23:00 and 7.7 at 03:00 if they stay awake. Five consecutive 4-h nights lower $S_w$ from 11.2 (one such night) to 8.9.

To the prior the model adds a time-on-task term $\gamma T(t)$, where $T$ is the time in hours since the start of the current task episode; it is frozen during a break and reset once the break has lasted 10~min. The value $\gamma = 0.4$ KSS~h$^{-1}$ is an assumption; time on task raises sleepiness in driving studies~\cite{akerstedt2010reaction}, but the Three-Process Model was built for average sedentary workloads~\cite{akerstedt2004predictions}, so the term may partly double count. Its sensitivity is reported in Section~\ref{sec:simulation}.

\subsection{State and dynamics}
The estimate is $s(t) = \mathrm{clip}\big(p(t) + \gamma T(t) + d(t)\big)$, where $d$ is the person's offset from the prior. The filter state is
\begin{equation}
  x = [\,d,\ \alpha_{\mathrm{RT}},\ \alpha_{\mathrm{BD}},\ \alpha_{\mathrm{PC}}\,]^\top,
\end{equation}
where the anchors $\alpha$ are how far this person's median reaction time (ms), blink duration (ms) and PERCLOS (\%) sit from the population at the same sleepiness. All components are random walks with $x_0 = 0$ and $P_0 = \mathrm{diag}(\sigma_d^2, \sigma_{\alpha,j}^2)$; over $\Delta$ hours $P \leftarrow P + \mathrm{diag}(q_d, q_{\alpha,j})\Delta$. The offset drifts quickly ($\sqrt{q_d} = 0.3$ KSS~h$^{-1/2}$) and the anchors slowly.

\subsection{Observations}
A KSS check-in $y$ observes the offset directly,
\begin{equation}
  y - p(t) - \gamma T(t) = d + \varepsilon_K, \qquad \varepsilon_K \sim \mathcal{N}(0, \sigma_K^2).
\end{equation}
Each passive feature $f_j$ (median reaction time of a test, or one minute's mean blink duration or PERCLOS) has a forward measurement model,
\begin{equation}
  f_j = a_j + b_j\,\big(s - 3\big) + \alpha_j + e_j,
  \label{eq:forward}
\end{equation}
where $a_j$ is the population mean at KSS~3 and $b_j$ the within-person slope. Substituting $s = p + \gamma T + d$ gives a linear observation of the state with $h = [b_j, 0, \ldots, 1, \ldots, 0]$ (1 in the position of $\alpha_j$) and offset $a_j + b_j(p + \gamma T - 3)$. The noise variance is
\begin{equation}
  R_j = \kappa_j\,\sigma_{e,j}^2 + \big(\sigma_{b,j}\,(\hat s - 3)\big)^2,
\end{equation}
where $\kappa_j = (1+\rho_j)/(1-\rho_j)$ allows for the lag-1 autocorrelation $\rho_j$ of successive minutes and the second term for between-person spread in the slope, evaluated at the current estimate $\hat s$. A reaction test is used only if it lasted the full 3~min. A webcam minute is used only if face coverage is at least 0.6; blink duration only if the minute has at least three blinks; and PERCLOS only if the PERCLOS guard (Section~\ref{sec:system}) allows it.

\paragraph*{Fitting on DROZY} The forward models were fitted on DROZY~\cite{massoz2016ulg} (Section~\ref{sec:drozy}) with subject fixed effects, so that the slope $b_j$ is a within-person slope and differences in people's baselines do not leak into it. The anchor's prior SD $\sigma_{\alpha,j}$ is the between-person SD of the baselines after removing their estimation noise. For the webcam features, the per-minute residual variance was split into a within-session part, used as $\sigma_{e,j}^2$, and a between-session part, modelled as drift of the anchor, at a rate that lets the anchor wander by the between-session residual SD over about 8~h. The slope spread $\sigma_{b,j}$ is the raw SD of per-person slopes from the subjects with three sessions, a conservative upper bound. No parameter was tuned on simulated people.

\paragraph*{Update} Each observation is absorbed with the standard Kalman update~\cite{kalman1960new},
\begin{equation}
  k = \frac{P h}{h^\top P h + R}, \quad x \leftarrow x + k\,(z - h^\top x), \quad P \leftarrow P - k\,h^\top P.
  \label{eq:kf}
\end{equation}

\subsection{What the anchors do}
A check-in pins $d$. A passive reading that then disagrees with the pinned level is explained mainly by its own anchor, because the anchor's prior variance is large compared with the remaining variance of $d$. Later readings from that signal move $d$ only to the extent that they change relative to the learned anchor. Check-ins therefore set the level and passive signals track change. Before any check-in, the anchor and the offset share each passive update according to their prior variances, so passive signals still move the estimate, but cautiously. A variant without anchors, used for comparison, fixes all $\alpha_j = 0$ and adds $\sigma_{\alpha,j}^2$ to $R_j$ as extra noise.

This design makes the person's ratings the reference. If a person's reaction times slow while their ratings stay flat, the model attributes the slowing to the anchor, which is the right behaviour for an estimate of how sleepy the person feels but hides objective decline~\cite{vandongen2003cumulative,leproult2003individual}. The application therefore also shows an \emph{objective check}: the latest median reaction time read through the population forward model alone, $\hat s_{\mathrm{obj}} = 3 + (f_{\mathrm{RT}} - a_{\mathrm{RT}})/b_{\mathrm{RT}}$, labelled as not adjusted to the person, with a note that a reading well above the person's own rating may mean they are under-rating their impairment.

\subsection{Forecast, what-if and attribution}
To forecast, the model holds the offset at its current mean and lets its variance grow, $\mathrm{sd}(t+\Delta) = (P_{dd} + q_d \Delta)^{1/2}$; the probability that latent sleepiness is at least 7 (``sleepy, but no effort to keep awake'') is $1 - \Phi\big((7 - \hat s)/\mathrm{sd}\big)$, and the predicted time of reaching KSS~7 is the first future time with $\hat s \geq 7$. The default projection assumes the person's usual rhythm of 50-min work blocks and 10-min breaks. What-if scenarios change the inputs: a 15-min break now, then the usual rhythm; a bedtime, which replaces the sleep episode used by the prior after waking; or working on without breaks, which lets $T$ grow. The interface also splits the offset into the sum of the updates each source has made to it (check-ins, reaction tests, webcam). This split is an accounting identity of the linear updates and depends on their order, so it is shown as a description of what the filter did, not as a causal attribution.

\begin{table}[t]
  \centering
  \caption{Model parameters in the released code}
  \label{tab:params}
  \footnotesize
  \setlength{\tabcolsep}{3pt}
  \begin{tabular}{@{}L{1.55cm}L{2.45cm}L{4.4cm}@{}}
    \toprule
    Symbol & Value & Meaning and source \\
    \midrule
    \multicolumn{3}{@{}l}{\emph{Three-Process Model}~\cite{akerstedt2004predictions,akerstedt2008accounting,ingre2014validating}} \\
    $L$, $H$ & 2.4, 14.3 & Asymptotes of $S$ \\
    $\delta_w$, $\delta_s$ & 0.0353, 0.381 h$^{-1}$ & Fall of $S$ awake; recovery asleep \\
    $B$ & 12.2 & Breakpoint; linear recovery below \\
    $M$, $\phi$ & 2.5; 16.8 h & Circadian amplitude, acrophase \\
    $U_a$, $U_m$, $\phi_U$ & 0.5, $-0.5$; 13.8 h & Ultradian amplitude, mesor, dip \\
    $W_0$, $\delta_i$ & $-5.72$; 1.51 h$^{-1}$ & Sleep inertia and its decay \\
    $a_K$, $b_K$ & 10.6, 0.6 & Alertness to KSS \\
    \midrule
    \multicolumn{3}{@{}l}{\emph{Offset and check-ins (assumed)}} \\
    $\gamma$ & 0.4 KSS h$^{-1}$ & Time on task; reset after a 10-min break \\
    $\sigma_d$, $\sqrt{q_d}$ & 1.2; 0.3 h$^{-1/2}$ & Prior SD and drift of the offset \\
    $\sigma_K$ & 0.75 & Noise of one KSS rating \\
    \midrule
    \multicolumn{3}{@{}l}{\emph{Forward models (DROZY, within person)}} \\
    Reaction test & $a$ 339.2 ms, $b$ 18.85 ms, $\sigma_e$ 36.4 & PVT median RT; anchor SD 48.7, drift 2 ms~h$^{-1/2}$, slope SD 7.9 \\
    Blink duration & $a$ 189.7 ms, $b$ 11.1 ms, $\sigma_e$ 24.6, $\kappa$ 1.04 & One minute, $\geq$3 blinks; anchor SD 22.9, drift 15 ms~h$^{-1/2}$, slope SD 14.3 \\
    PERCLOS & $a$ 3.24\%, $b$ 0.62\%, $\sigma_e$ 2.4, $\kappa$ 1.93 & One minute, PERCLOS guard; anchor SD 2.38, drift 0.83\%~h$^{-1/2}$, slope SD 0.47 \\
    \bottomrule
  \end{tabular}
\end{table}
